\section{Shortest Path in an Undirected Graph with Unit Weights}
\label{sec:graph-shortest-path-unit-weights}

\problemstatement{Given an undirected graph where every edge has
weight $1$, and a source node, find the shortest distance from the
source to every other node.}

\begin{patternbox}
\emph{``Shortest path, every edge the same cost''} is exactly what
plain BFS already computes: BFS processes nodes in strictly
non-decreasing order of distance from the source, so the level at which
a node is first reached \textbf{is} its shortest distance. No
relaxation, priority queue, or weight comparison is needed at all.
\end{patternbox}

\subsection*{Why unit weights make BFS sufficient}

The general shortest-path relaxation rule,
$\textit{dist}[v] = \min(\textit{dist}[v], \textit{dist}[u]+w)$,
collapses to $\textit{dist}[v] = \textit{dist}[u]+1$ when every $w=1$ --
and BFS's queue already guarantees every node is dequeued with its
\emph{final}, minimal distance already correctly assigned (this is the
same property exploited by multi-source BFS in the Graph Basics
chapter's 01 Matrix problem, here with a single source instead of
many).

\subsection*{Algorithm}

\begin{enumerate}
  \item Initialize $\textit{dist}[\textit{src}]=0$ and every other
        $\textit{dist}[i]=\infty$.
  \item BFS from the source: for each neighbor of the current node
        with $\textit{dist}[\textit{neighbor}] > \textit{dist}[\textit{node}]+1$,
        update it and push the neighbor.
\end{enumerate}

\subsection*{C++ implementation}

\begin{lstlisting}
#include <bits/stdc++.h>
using namespace std;

vector<int> shortestPathUnitWeights(int n, vector<vector<int>>& adj, int src) {
    vector<int> dist(n, INT_MAX);
    dist[src] = 0;
    queue<int> q;
    q.push(src);

    while (!q.empty()) {
        int node = q.front();
        q.pop();

        for (int neighbor : adj[node]) {
            if (dist[node] + 1 < dist[neighbor]) {
                dist[neighbor] = dist[node] + 1;
                q.push(neighbor);
            }
        }
    }
    return dist;
}

int main() {
    int n = 5;
    vector<vector<int>> adj(n);
    auto addEdge = [&](int u, int v) { adj[u].push_back(v); adj[v].push_back(u); };
    addEdge(0, 1); addEdge(0, 2); addEdge(1, 3); addEdge(2, 3); addEdge(3, 4);

    vector<int> dist = shortestPathUnitWeights(n, adj, 0);
    for (int d : dist) cout << d << " ";
    cout << endl; // 0 1 1 2 3
    return 0;
}
\end{lstlisting}

\subsection*{Dry run}

\dryrun{Take the graph above, source $0$.}

$\textit{dist}=[0,\infty,\infty,\infty,\infty]$. Pop $0$: relax $1$
($\infty\to1$) and $2$ ($\infty\to1$); push both. Pop $1$: relax $3$
($\infty\to2$); push. Pop $2$: neighbor $3$ already at distance $2$,
and $\textit{dist}[2]+1=2$ is not strictly smaller, no update. Pop
$3$: relax $4$ ($\infty\to3$); push. Pop $4$: no unrelaxed neighbors.
Final: $[0,1,1,2,3]$.

\begin{complexitybox}
\textbf{Time Complexity.} $O(V+E)$.
\textbf{Space Complexity.} $O(V)$ for \textit{dist} and the queue.
\end{complexitybox}

\begin{takeawaybox}
This problem is included specifically as a baseline: every other
algorithm in this chapter is, at its core, solving the same
``relax edges until distances stabilize'' problem that unit-weight BFS
solves trivially -- they exist only because weighted, negative, or
all-pairs distances break BFS's clean level-by-level guarantee in
different specific ways.
\end{takeawaybox}
